Pain presentations in human myocarditis and any links or distinctions involving aortic pain.
Human studies on myocarditis and aortic (or “aortic-like”) pain primarily highlight differential diagnosis challenges, rare overlapping presentations, and secondary effects of inflammation on aortic properties rather than a direct causal link between myocarditis and classic aortic pain.
Myocarditis (inflammation of the heart muscle) most commonly presents with chest pain, while true aortic pathology (dissection, aortitis, aneurysm) produces distinct pain patterns. Studies and case reports show occasional mimicry or co-occurrence.
Typical Pain Characteristics
Myocarditis: Chest pain is the dominant symptom in 85–95% of acute cases (often 93% in large cohorts). It is frequently infarct-like (pressure, tightness) or pleuritic (sharp, positional, worse with inspiration or lying flat if pericardium is involved). Associated features include recent viral prodrome (fever, respiratory/gastrointestinal symptoms in up to 80%), elevated troponin, ECG changes (ST elevation or T-wave inversion), and dyspnea or arrhythmias. Pain may radiate but is less commonly described as “tearing” or primarily interscapular.
Aortic dissection / acute aortic syndromes: Sudden-onset severe “tearing” or “ripping” pain, maximal at onset, often radiating to the back, abdomen, or neck. Pulse deficits, blood-pressure asymmetry, or neurological signs raise suspicion. Classic aortic pain is less positional than pericarditic pain.
Key Human Studies and Case Reports on Overlap or Mimicry
Myocarditis mimicking aortic dissection
A well-documented 2019 case (BMJ Case Reports) described a 28-year-old woman with sudden-onset severe upper back pain, mild chest discomfort, and bilateral arm paresthesia. Prehospital ECG showed inferior ST-segment elevation with reciprocal changes, raising concern for both inferior myocardial infarction and ascending aortic dissection. CT aortogram excluded dissection; coronary angiography showed unobstructed arteries; cardiac MRI confirmed acute myocarditis (edema and mid-wall/epicardial late gadolinium enhancement in multiple segments). Troponin peaked at ~3000 ng/L; the patient recovered with conservative therapy (beta-blocker, ACE inhibitor). This illustrates how myocarditis can produce referred “aortic-like” back pain and ECG changes that force urgent exclusion of life-threatening aortic or coronary disease.
Aortic stiffness in acute myocarditis
A CMR study compared pulse-wave velocity (PWV, a marker of aortic stiffness) in patients with acute myocarditis (n=44), chronic systemic inflammatory diseases (RA/SLE, n=55), and healthy controls (n=38). Both inflammatory groups showed significantly elevated central PWV versus controls (myocarditis ~8.4 m/s, chronic inflammation ~8.5 m/s, controls ~5.1 m/s). The increase in acute (self-limited) myocarditis was comparable in magnitude to sustained inflammation, indicating that even short-lived systemic inflammation can acutely stiffen the aorta. Age correlated with PWV in all groups.
A separate pediatric/young-adult follow-up study found no persistent differences in aortic or carotid intima-media thickness, PWV, or distensibility years after recovered viral myocarditis compared with controls, suggesting the stiffness change is largely reversible once inflammation resolves.
Myocarditis + aortitis / large-vessel vasculitis  
Myocarditis is a rare (~2–3%) but recognized presentation of Takayasu arteritis or giant-cell arteritis (large-vessel vasculitis that frequently involves the aorta). Patients may present with chest pain, heart failure, or elevated biomarkers alongside aortic-wall thickening, stenosis, or aneurysm. Imaging (CMR, PET-CT, CTA) and immunosuppression are key.
A post-COVID case report described simultaneous autoimmune myocarditis (EF drop to 30%, high anticardiac antibodies, MRI-confirmed) and aortitis with progressive infrarenal aortic aneurysm developing one month after SARS-CoV-2 infection. Corticosteroids improved the myocarditis, but the aneurysm required surgery; viral proteins were detected in inflammatory infiltrates.
Rare reports link fulminant myocarditis with aortic thrombosis or occlusion in the setting of severe systemic inflammation and cytokine storm.
Diagnostic Implications from Clinical Series
Large registries (e.g., MyocarditIRM, Korean multicenter data, ACC 2024 expert consensus) emphasize three classic myocarditis presentations: chest-pain phenotype, heart-failure/shock phenotype, and arrhythmia phenotype. Chest pain is by far the most common. Because ECG and biomarker changes overlap with acute coronary syndromes and occasionally aortic syndromes, guidelines stress rapid exclusion of dissection (CT aortogram) and coronary occlusion (angiography or CT) before attributing findings to myocarditis. Cardiac MRI (Lake Louise criteria) is the preferred non-invasive confirmatory test in stable patients; endomyocardial biopsy is reserved for high-risk or unclear cases.
Summary of Evidence Strength
Direct “aortic pain” caused by pure myocarditis is uncommon; the strongest evidence is the case-level demonstration of mimicry (back pain + ECG changes prompting dissection work-up).
Inflammation from myocarditis can acutely increase aortic stiffness (PWV data), providing a plausible biomechanical link to altered aortic wall properties, though this does not typically produce dissection-level pain.
True concurrent aortitis and myocarditis occurs mainly in the setting of systemic vasculitis or post-infectious autoimmunity and requires specific imaging and immunosuppressive strategies.
These studies underscore the need for broad differential diagnosis when chest or back pain is accompanied by myocardial injury markers, especially in younger patients or those with recent infection. Imaging (CT aortogram + coronary assessment + CMR) remains essential to separate myocarditis from primary aortic pathology.
